%% 由 build/emit_tex.py 生成（生成物，勿手改）；定制请放 user_style.tex / user_meta.yaml
\chapter{图论与网络}\label{ux56feux8bbaux4e0eux7f51ux7edc}

\begin{quote}
本附录为第 6 章 NetworkX
提供图论与算法背景：图表示、基本量、经典算法（BFS/最短路/MST/中心性/社区），并补上本科常见难点（握手定理、树的性质、模块度）。
\end{quote}

\begin{center}\rule{0.5\linewidth}{0.5pt}\end{center}

\section{D.1
直觉故事：朋友圈里的''谁最重要''}\label{d.1-ux76f4ux89c9ux6545ux4e8bux670bux53cbux5708ux91ccux7684ux8c01ux6700ux91cdux8981}

``谁最重要''没有唯一答案：认识 10 个互相不认识的人 → 桥梁；认识 10
个彼此都认识的人 → 不是桥梁。所以：

\begin{itemize}
\tightlist
\item
  \textbf{度中心性}：连接数多 → ``活跃''；
\item
  \textbf{接近中心性}：到所有人平均距离短 → ``信息传得快''；
\item
  \textbf{介数中心性}：很多最短路经过它 → ``枢纽/瓶颈''；
\item
  \textbf{PageRank}：被''重要的人''指向 → ``影响力''。
\end{itemize}

做网络分析第一步是\textbf{明确要问什么}，再选指标。

\begin{quote}
\textbf{正文见}：\href{../../chapters/06-networkx/02-分析图.md}{6
NetworkX · 02
分析图}、\href{../../chapters/06-networkx/03-求解图的基本问题.md}{6
NetworkX · 03 求解图的基本问题}。
\end{quote}

\begin{center}\rule{0.5\linewidth}{0.5pt}\end{center}

\section{D.2 图的基本表示（讲解 +
手算）}\label{d.2-ux56feux7684ux57faux672cux8868ux793aux8bb2ux89e3-ux624bux7b97}

\subsection{D.2.1
怎么把图''喂给电脑''}\label{d.2.1-ux600eux4e48ux628aux56feux5582ux7ed9ux7535ux8111}

{\def\LTcaptype{none} % do not increment counter
\begin{longtable}[]{@{}lll@{}}
\toprule\noalign{}
表示 & 样子 & 适合 \\
\midrule\noalign{}
\endhead
\bottomrule\noalign{}
\endlastfoot
邻接矩阵 \(A\) & \(A[i][j]=1\)（或权重） &
稠密图、矩阵运算；点多费内存 \\
邻接表 & 每个点存邻居列表 & 稀疏大图；遍历快 \\
\end{longtable}
}

手算例子（4 点：A-B、A-C、B-C、C-D）：

\[A=\begin{bmatrix}0&1&1&0\\1&0&1&0\\1&1&0&1\\0&0&1&0\end{bmatrix}\]

（无向图对称；\texttt{np.linalg.eigvalsh(A)}
能直接读图的性质，如连通性/谱聚类。）

\subsection{D.2.2
基本量：度、连通、树}\label{d.2.2-ux57faux672cux91cfux5ea6ux8fdeux901aux6811}

\begin{itemize}
\tightlist
\item
  \textbf{度} \(\deg(v)\)：邻边数；有向图分\textbf{入度/出度}；
\item
  \textbf{握手定理}：\(\sum_v \deg(v)=2|E|\)（每条边贡献 2
  度）------检查数据是否''边数对着''；
\item
  \textbf{路径/连通分量}：可达性；\textbf{树}：无环连通图，\(n\) 个点恰
  \(n-1\) 条边；
\item
  \textbf{二分图}：点分两类，边只在类间（学生-选课、用户-商品）。
\end{itemize}

\subsection{D.2.3
加权与有向}\label{d.2.3-ux52a0ux6743ux4e0eux6709ux5411}

\begin{itemize}
\tightlist
\item
  \textbf{加权图}：边带距离/成本/强度；最短路、MST、网络流依赖这些权；
\item
  \textbf{有向图}：边有方向（关注、引用、资金流）；PageRank/拓扑排序只在有向图上。
\end{itemize}

\begin{center}\rule{0.5\linewidth}{0.5pt}\end{center}

\section{D.3
经典算法：思路与手算}\label{d.3-ux7ecfux5178ux7b97ux6cd5ux601dux8defux4e0eux624bux7b97}

\subsection{D.3.1 遍历：BFS 与
DFS}\label{d.3.1-ux904dux5386bfs-ux4e0e-dfs}

\begin{itemize}
\tightlist
\item
  \textbf{BFS}：一层层扩散。无权图最短路靠它；手算：从 A
  出发，A→B、C，再 C→D，路径 A→C→D 两步最短；
\item
  \textbf{DFS}：一条路走到黑再回头；用于检测环、连通分量、拓扑排序。
\end{itemize}

\subsection{D.3.2 最短路径}\label{d.3.2-ux6700ux77edux8defux5f84}

{\def\LTcaptype{none} % do not increment counter
\begin{longtable}[]{@{}llll@{}}
\toprule\noalign{}
算法 & 适用 & 思路 & 复杂度 \\
\midrule\noalign{}
\endhead
\bottomrule\noalign{}
\endlastfoot
BFS & 无权 & 按层扩散 & \(O(V+E)\) \\
Dijkstra & 非负权 & 每次取当前最近点并松弛 & \(O((V+E)\log V)\)（堆） \\
Bellman-Ford & 可能有负权 & 反复松弛 & \(O(VE)\) \\
Floyd-Warshall & 全源 & 动态规划 & \(O(V^3)\) \\
\end{longtable}
}

手算 Dijkstra（A-B=2、A-C=5、B-C=1、C-D=4）：A=0 → B=2、C=5 → 取
B=2，松弛 C 得 3 → 取 C=3，松弛 D=7 → 得 A→B→C→D，距离 7。

\subsection{D.3.3
最小生成树（MST）}\label{d.3.3-ux6700ux5c0fux751fux6210ux6811mst}

目标：总权重最小的连通子图（树）。\textbf{Kruskal}：边按权排序，用并查集''不成环就加''；最后正好
\(n-1\)
条边。\textbf{Prim}：从一个点出发，每次加入''连接已选集合的最便宜边''。两者都满足''贪心
+ 无环''的树性质。

\subsection{D.3.4
中心性与社区}\label{d.3.4-ux4e2dux5fc3ux6027ux4e0eux793eux533a}

\begin{itemize}
\tightlist
\item
  \textbf{介数中心性}：统计每条最短路经过某点的次数；
\item
  \textbf{PageRank}：随机游走稳态；链接不仅''数量''还看''质量''；
\item
  \textbf{社区发现}：把网络切成''内部紧密、外部稀疏''。\textbf{模块度}：
\end{itemize}

\[Q=\frac{1}{2m}\sum_{ij}\Big[A_{ij}-\frac{d_i d_j}{2m}\Big]\delta(c_i,c_j)\]

直觉：实际边数 - 随机网络下的期望边数（同一社区内）。\textbf{Louvain}
贪心优化模块度，快且常用；结果与分辨率/随机种子有关，\textbf{要报告方法与参数}。

\begin{center}\rule{0.5\linewidth}{0.5pt}\end{center}

\section{D.4
动手例题（选做）}\label{d.4-ux52a8ux624bux4f8bux9898ux9009ux505a}

\textbf{例：Dijkstra 代码验证}：

\begin{Shaded}
\begin{Highlighting}[]
\ImportTok{import}\NormalTok{ networkx }\ImportTok{as}\NormalTok{ nx}
\NormalTok{G }\OperatorTok{=}\NormalTok{ nx.Graph()}
\NormalTok{G.add\_weighted\_edges\_from([(}\StringTok{"A"}\NormalTok{,}\StringTok{"B"}\NormalTok{,}\DecValTok{2}\NormalTok{), (}\StringTok{"A"}\NormalTok{,}\StringTok{"C"}\NormalTok{,}\DecValTok{5}\NormalTok{), (}\StringTok{"B"}\NormalTok{,}\StringTok{"C"}\NormalTok{,}\DecValTok{1}\NormalTok{), (}\StringTok{"C"}\NormalTok{,}\StringTok{"D"}\NormalTok{,}\DecValTok{4}\NormalTok{)])}
\BuiltInTok{print}\NormalTok{(nx.shortest\_path(G, }\StringTok{"A"}\NormalTok{, }\StringTok{"D"}\NormalTok{, weight}\OperatorTok{=}\StringTok{"weight"}\NormalTok{))          }\CommentTok{\# [\textquotesingle{}A\textquotesingle{},\textquotesingle{}B\textquotesingle{},\textquotesingle{}C\textquotesingle{},\textquotesingle{}D\textquotesingle{}]}
\BuiltInTok{print}\NormalTok{(nx.shortest\_path\_length(G, }\StringTok{"A"}\NormalTok{, }\StringTok{"D"}\NormalTok{, weight}\OperatorTok{=}\StringTok{"weight"}\NormalTok{))   }\CommentTok{\# 7}
\end{Highlighting}
\end{Shaded}

\textbf{例：中心性完全不同}。星型网络（中心 A 连 5 个叶子）：A
的介数中心性远高于叶子；但每个叶子的度都是
1------``度中心性''看不出谁是枢纽。

\begin{center}\rule{0.5\linewidth}{0.5pt}\end{center}

\section{D.5 Python
对应（速查）}\label{d.5-python-ux5bf9ux5e94ux901fux67e5}

\begin{Shaded}
\begin{Highlighting}[]
\ImportTok{import}\NormalTok{ networkx }\ImportTok{as}\NormalTok{ nx}
\NormalTok{G }\OperatorTok{=}\NormalTok{ nx.Graph()}
\NormalTok{G.add\_edges\_from([(}\StringTok{"A"}\NormalTok{,}\StringTok{"B"}\NormalTok{,\{}\StringTok{"w"}\NormalTok{:}\DecValTok{2}\NormalTok{\}), (}\StringTok{"A"}\NormalTok{,}\StringTok{"C"}\NormalTok{,\{}\StringTok{"w"}\NormalTok{:}\DecValTok{5}\NormalTok{\}),}
\NormalTok{                  (}\StringTok{"B"}\NormalTok{,}\StringTok{"C"}\NormalTok{,\{}\StringTok{"w"}\NormalTok{:}\DecValTok{1}\NormalTok{\}), (}\StringTok{"C"}\NormalTok{,}\StringTok{"D"}\NormalTok{,\{}\StringTok{"w"}\NormalTok{:}\DecValTok{4}\NormalTok{\})])}

\NormalTok{nx.shortest\_path(G, }\StringTok{"A"}\NormalTok{, }\StringTok{"D"}\NormalTok{, weight}\OperatorTok{=}\StringTok{"w"}\NormalTok{)      }\CommentTok{\# Dijkstra}
\NormalTok{nx.shortest\_path\_length(G, }\StringTok{"A"}\NormalTok{, }\StringTok{"D"}\NormalTok{, weight}\OperatorTok{=}\StringTok{"w"}\NormalTok{)}
\NormalTok{nx.minimum\_spanning\_tree(G, weight}\OperatorTok{=}\StringTok{"w"}\NormalTok{)}
\NormalTok{nx.degree\_centrality(G)}\OperatorTok{;}\NormalTok{ nx.betweenness\_centrality(G)}
\NormalTok{nx.pagerank(G)}
\NormalTok{nx.community.louvain\_communities(G)}
\NormalTok{nx.draw(G, with\_labels}\OperatorTok{=}\VariableTok{True}\NormalTok{)}
\end{Highlighting}
\end{Shaded}

{\def\LTcaptype{none} % do not increment counter
\begin{longtable}[]{@{}ll@{}}
\toprule\noalign{}
你想要 & 用哪个 \\
\midrule\noalign{}
\endhead
\bottomrule\noalign{}
\endlastfoot
最短路 & \texttt{nx.shortest\_path}（加权传 \texttt{weight}） \\
MST & \texttt{nx.minimum\_spanning\_tree} \\
中心性 & \texttt{degree\ /\ closeness\ /\ betweenness\ /\ pagerank} \\
社区 & \texttt{nx.community} \\
可视化 & \texttt{nx.draw}（或结合 Matplotlib） \\
\end{longtable}
}

\begin{center}\rule{0.5\linewidth}{0.5pt}\end{center}

\section{D.6 常见误区}\label{d.6-ux5e38ux89c1ux8befux533a}

{\def\LTcaptype{none} % do not increment counter
\begin{longtable}[]{@{}ll@{}}
\toprule\noalign{}
误区 & 正确 \\
\midrule\noalign{}
\endhead
\bottomrule\noalign{}
\endlastfoot
以为图只能是邻接矩阵 & 稀疏大图用 NetworkX 的邻接表实现 \\
无向图当成有向图算 & \texttt{Graph} 与 \texttt{DiGraph} 分开 \\
用最短路回答''谁最重要'' & 先明确''重要''的定义 \\
忽略权重 & 加权算法要传 \texttt{weight} \\
社区数量随便定 & 社区发现与分辨率有关；试多组参数并看模块度 \\
握手定理没验证 & 先看 \texttt{sum\ deg(v)\ ==\ 2*\textbar{}E\textbar{}}
排查数据 \\
\end{longtable}
}

\begin{center}\rule{0.5\linewidth}{0.5pt}\end{center}

\section{D.7
使用章节（双向）}\label{d.7-ux4f7fux7528ux7ae0ux8282ux53ccux5411}

{\def\LTcaptype{none} % do not increment counter
\begin{longtable}[]{@{}
  >{\raggedright\arraybackslash}p{(\linewidth - 4\tabcolsep) * \real{0.3333}}
  >{\raggedright\arraybackslash}p{(\linewidth - 4\tabcolsep) * \real{0.3333}}
  >{\raggedright\arraybackslash}p{(\linewidth - 4\tabcolsep) * \real{0.3333}}@{}}
\toprule\noalign{}
\begin{minipage}[b]{\linewidth}\raggedright
章
\end{minipage} & \begin{minipage}[b]{\linewidth}\raggedright
哪里用到
\end{minipage} & \begin{minipage}[b]{\linewidth}\raggedright
链接
\end{minipage} \\
\midrule\noalign{}
\endhead
\bottomrule\noalign{}
\endlastfoot
6 NetworkX & 建图/分析/最短路 &
\href{../../chapters/06-networkx/01-创建图.md}{01 创建图} \\
6 NetworkX & 中心性/社区 &
\href{../../chapters/06-networkx/02-分析图.md}{02 分析图} \\
6 NetworkX & 综合案例 &
\href{../../chapters/06-networkx/04-综合案例.md}{04 综合案例} \\
\end{longtable}
}

\textbf{下游衔接}：intro-mathmodel 第 4 章（复杂网络与图论模型）。
\textbf{延伸阅读}：NetworkX
官方教程、图论教材、\href{./references.md}{references.md}。

\begin{center}\rule{0.5\linewidth}{0.5pt}\end{center}

\section{D.10 常见考题与自查（考前 10
分钟）}\label{d.10-ux5e38ux89c1ux8003ux9898ux4e0eux81eaux67e5ux8003ux524d-10-ux5206ux949f}

{\def\LTcaptype{none} % do not increment counter
\begin{longtable}[]{@{}lll@{}}
\toprule\noalign{}
会了吗？ & 考点 & 一句话答案 \\
\midrule\noalign{}
\endhead
\bottomrule\noalign{}
\endlastfoot
□ & BFS 与 DFS 适用 & BFS 无权最短路；DFS 环/连通/拓扑 \\
□ & Dijkstra 前提 & 边权非负 \\
□ & MST 贪心 & Kruskal/Prim 每次选不形成环的最便宜边 \\
□ & 度 vs 介数中心性 & 度=活跃；介数=桥梁/瓶颈 \\
□ & 模块度 & 社区内实际边 - 随机期望，越大越''像社区'' \\
□ & 握手定理 & 度之和 = 2×边数 \\
\end{longtable}
}

\begin{center}\rule{0.5\linewidth}{0.5pt}\end{center}

\section{D.9
综合案例：给一个小网络做''体检''}\label{d.9-ux7efcux5408ux6848ux4f8bux7ed9ux4e00ux4e2aux5c0fux7f51ux7edcux505aux4f53ux68c0}

\textbf{问题}：给定学校 20
人社交网络，找出''信息传播最快的人''、``最容易被绕过的人''与''天然小团伙''。

步骤：

\begin{enumerate}
\def\labelenumi{\arabic{enumi}.}
\tightlist
\item
  建图：读边表，检查握手定理（度之和 = 2×边数）；
\item
  最短路：无权图用 BFS，加权用 Dijkstra；
\item
  中心性：同时看度/接近/介数/PageRank，画四张图对比（结论不同是正常的，要把定义说清）；
\item
  社区：Louvain 跑几组分辨率，报告模块度与参数；
\item
  输出：``若做疫情通报，优先通知介数中心性最高的人；若做班委选举，关注度中心性与
  PageRank 都高的人。''
\end{enumerate}

参考要点代码：

\begin{verbatim}
import networkx as nx
G = nx.read_edgelist("edges.txt")
print(sum(d for _, d in G.degree()) == 2 * G.number_of_edges())
print(nx.betweenness_centrality(G)); print(nx.pagerank(G))
from networkx.algorithms.community import louvain_communities
print(louvain_communities(G, seed=0))
\end{verbatim}

\textbf{反思}：图分析''结论''依赖指标定义；把''问什么''和''用什么''一起写进报告。

\begin{center}\rule{0.5\linewidth}{0.5pt}\end{center}

\section{D.8
例题集（深入练习）}\label{d.8-ux4f8bux9898ux96c6ux6df1ux5165ux7ec3ux4e60}

\textbf{例 1：Kruskal 手算}。边（A-B=2, A-C=5, B-C=1, C-D=4）：排序
B-C(1) → A-B(2) → C-D(4) →
A-C(5)。依次加入不成环的边：B-C、A-B、C-D；此时 4 点 3 边成树，总权
7。A-C 会成环，跳过。

\textbf{例 2：PageRank 三点手算}（有向环 A→B→C→A
加自环简化不计）：稳态分布与转移矩阵最大特征向量相同------直观上''每个点的权重等于流入它的点的权重之和''。

\textbf{例 3：BFS 手算}。图 A-B、A-C、B-D、C-D：BFS 从 A 出发的层序：0
层 A；1 层 B、C；2 层 D。因此 A 到 D 最短路为 2（经 B 或
C），无权图最短路 = 层数。
